\section{Predictive usefulness and computational cost}
\label{sec:utility-question}
Let $X\in\R^d$ be an observed history, $Y\in\R^{d_y}$ its future
response, and $\Gamma_{\mathcal D}$ a representation fitted on development
data $\mathcal D$. Expectations below are conditional on the fitted
state. Each experiment specifies its own evaluation distribution and
loss aggregation. A deterministic representation can discard predictive
information; appending it to $X$ preserves the raw history and hence its
Bayes risk (Proposition~\ref{prop:nf-raw-preservation}). This does not
ensure equal performance of finite, regularized fitted decoders.

\paragraph{Residual information relative to a baseline.}
Fix $f_0:\R^d\to\R^{d_y}$, put $e=Y-f_0(X)$ and
$Z=\Gamma_{\mathcal D}(X)$, and assume $\E\|e\|_2^2<\infty$.
Write $R(f)=\E\|Y-f(X)\|_2^2$ and let
$L^2(P_Z;\R^{d_y})$ contain the measurable corrections with finite
second moment under the distribution $P_Z$ of $Z$. Conditional
expectation gives the best correction $q^*(Z)=\E[e\mid Z]$ and
\begin{equation}\label{eq:utility-oracle}
 R(f_0)-\inf_{q\in L^2(P_Z;\R^{d_y})}
 R(f_0+q\circ\Gamma_{\mathcal D})
 =\E\|\E[e\mid Z]\|_2^2.
\end{equation}
The circle denotes composition. Appendix~\ref{app:utility-projection}
proves the identity. Improvement is possible in this unrestricted
correction class precisely when the residual conditional mean is
nonzero with positive probability. The identity does not determine the
performance of a correction estimated in a finite class.

\paragraph{Fitted procedures and computation.}
For fitted forecasts $\widehat f_0$ and $\widehat f_s$, evaluation unit
$\mathcal E$, and specified squared loss $\ell$, define
\begin{equation}\label{eq:utility-delta}
 \Delta_s=\E[\ell(\widehat f_0;\mathcal E)
             -\ell(\widehat f_s;\mathcal E)\mid\mathcal D].
\end{equation}
Positive $\Delta_s$ favors the richer procedure. A comparison may require
a minimum gain $\delta>0$ and cost budget $B>0$, with measured complete
procedure cost $C_s\leq B$. Feature construction, fitting, and selection
are charged even when the selected prediction equals the baseline.
Evaluation outcomes cannot enter fitting or selection. Independently
sampled repetitions, overlapping windows, and selected source families
supply different units of evidence. Missing predictions without a
prespecified fallback make the affected full-roster aggregate undefined;
restriction to successful fits changes the evaluated population. No
experiment here estimates unrestricted natural-data Bayes risk or
certifies the minimum complexity necessary to achieve a target accuracy.

\paragraph{A correction that retains the fitted baseline.}
For geometric family $g$, fitted correction $\widehat q_g$, and scalar
mixing weight $\gamma$, define
\begin{equation}\label{eq:utility-additive}
 \widehat f_{g,\gamma}(x)
 =\widehat f_0(x)+\gamma\widehat q_g(x),\qquad 0\leq\gamma\leq1.
\end{equation}
Zero correction recovers the same fitted baseline. Residual fitting is
established \citep{Buhlmann2003L2}; these studies compare the geometric
features supplied to the correction. With square-integrable residual
and correction, expansion of squared loss yields
\begin{equation}\label{eq:utility-alignment}
 R(\widehat f_{g,\gamma})-R(\widehat f_0)
 =\gamma^2\E\|\widehat q_g(X)\|_2^2
 -2\gamma\E\langle Y-\widehat f_0(X),\widehat q_g(X)\rangle.
\end{equation}
Here $\langle\cdot,\cdot\rangle$ is the Euclidean inner product.
For a correction nonzero with positive probability, positive expected
alignment is necessary and sufficient for improvement at some
sufficiently small positive $\gamma$. A finite candidate grid may miss
that interval, and validation can select a worse forecast. Baseline
availability is therefore a candidate-class property, not a performance
guarantee; Appendix~\ref{app:utility-alignment} gives the derivation.
